% SEGMENT OLP-0110-S01
% Part: first-order-logic
% Chapter: tableaux
% Section: identity

\documentclass[../../../include/open-logic-section]{subfiles}

\begin{document}

\olfileid{fol}{tab}{ide}

\olsection{\usetoken{S}{identity} கொண்ட \usetoken{P}{tableau}}

!!^{identity} கொண்ட !!^{tableau}s க்குக் கூடுதல் அனுமான விதிகள் தேவை.
$\eq$ க்கான விதிகள் பின்வருமாறு ($t$, $t_1$, $t_2$ ஆகியவை மூடிய சொற்கள்):

\begin{defish}
\AxiomC{}
\RightLabel{$\eq$}
\UnaryInfC{\sFmla{\True}{\eq[t][t]}}
\DisplayProof
\hfill
\AxiomC{\sFmla{\True}{\eq[t_1][t_2]}}
\noLine
\UnaryInfC{\sFmla{\True}{!A(t_1)}}
\RightLabel{$\TRule{\True}{\eq}$}
\UnaryInfC{\sFmla{\True}{!A(t_2)}}
\DisplayProof
\hfill
\AxiomC{\sFmla{\True}{\eq[t_1][t_2]}}
\noLine
\UnaryInfC{\sFmla{\False}{!A(t_1)}}
\RightLabel{$\TRule{\False}{\eq}$}
\UnaryInfC{\sFmla{\False}{!A(t_2)}}
\DisplayProof
\end{defish}
மற்ற எல்லா விதிகளுக்கும் மாறாக, $\TRule{\True}{\eq}$,
$\TRule{\False}{\eq}$ ஆகியவற்றைப் பயன்படுத்தக் கிளையில் ஏற்கெனவே
\emph{இரண்டு} குறியிட்ட !!{formula}s இருக்க வேண்டும் என்பதைக் கவனிக்கவும்;
அவை $\sFmla{\True}{\eq[t_1][t_2]}$, $\sFmla{S}{!A(t_1)}$ ஆகிய இரண்டும்.

% SEGMENT OLP-0110-S02
\begin{ex}
$s$, $t$ ஆகியவை மூடிய சொற்கள் எனில், $\eq[s][t], !A(s)
\Proves !A(t)$:
\begin{oltableau}
  [\sFmla{\False}{\formula{A}(t)}, just = \TAss
    [\sFmla{\True}{\eq[s][t]}, just = \TAss
      [\sFmla{\True}{\formula{A}(s)}, just = \TAss
        [\sFmla{\True}{\formula{A}(t)}, just={\TRule{\True}{\eq}[2, 3]}, close]
      ]
    ]
  ]
\end{oltableau}
இது ஒன்றானவற்றின் பிரதியிடத்தக்கதன்மைக் கொள்கை அல்லது லைப்னிட்சின் விதி
என்று அறிமுகமானதாக இருக்கலாம்.

% SEGMENT OLP-0110-S03
!!^{tableau}s, $\eq$ சமச்சீரானது என்பதை, அதாவது $\eq[s_1][s_2]
\Proves \eq[s_2][s_1]$ என்பதை, நிறுவுகின்றன:
\begin{oltableau}
  [\sFmla{\False}{\eq[s_2][s_1]}, just = \TAss
    [\sFmla{\True}{\eq[s_1][s_2]}, just = \TAss
      [\sFmla{\True}{\eq[s_1][s_1]}, just = {$\eq$}
        [\sFmla{\True}{\eq[s_2][s_1]}, just = {\TRule{\True}{\eq}[2, 3]}, close]
      ]
    ]
  ]
\end{oltableau}
இங்கு வரி~$2$, $\TRule{\True}{\eq}$ க்குத் தேவையான முதல்
!!{formula}~$\sFmla{\True}{\eq[s_1][s_2]}$. வரி~$3$ அதற்குத் தேவையான
இரண்டாவது வாய்பாடு; அது $\sFmla{\True}{!A(s_1)}$ என்ற வடிவிலுள்ளது—
% SOURCE CORRECTION TA-TAB-007: the second premise is A(s_1), not A(s_2).
$!A(x)$ ஐ $\eq[x][s_1]$ எனக் கருதுக; அப்போது $!A(s_1)$ என்பது
$\eq[s_1][s_1]$, $!A(s_2)$ என்பது $\eq[s_2][s_1]$.

% SEGMENT OLP-0110-S04
$\eq$ கடப்புப் பண்புடையது என்பதையும், அதாவது $\eq[s_1][s_2],
\eq[s_2][s_3] \Proves \eq[s_1][s_3]$ என்பதையும், அவை நிறுவுகின்றன:
\begin{oltableau}
  [\sFmla{\False}{\eq[s_1][s_3]}, just = \TAss
    [\sFmla{\True}{\eq[s_1][s_2]}, just = \TAss
      [\sFmla{\True}{\eq[s_2][s_3]}, just = \TAss
        [\sFmla{\True}{\eq[s_1][s_3]}, just = {\TRule{\True}{\eq}[3, 2]}, close]
      ]
    ]
  ]
\end{oltableau}
இந்த !!{tableau}வில் $\TRule{\True}{\eq}$ க்குத் தேவையான முதல் !!{formula},
வரி~$3$ இலுள்ள $\sFmla{\True}{\eq[s_2][s_3]}$ ஆகும் ($s_2$ ஆனது~$t_1$
இன் பாத்திரத்தையும் $s_3$ ஆனது~$t_2$ இன் பாத்திரத்தையும் வகிக்கின்றன).
தேவையான இரண்டாவது வாய்பாடு வரி~$2$; அது
$\sFmla{\True}{!A(s_2)}$ என்ற வடிவிலுள்ளது. இங்கு $!A(x)$ ஐ
$\eq[s_1][x]$ எனக் கருதுக; அப்போது $!A(s_2)$ என்பது
% SOURCE CORRECTION TA-TAB-007: instantiate the displayed schema as eq[s_1][s_2].
$\eq[s_1][s_2]$ (அதாவது வரி~$2$) ஆகவும், $!A(s_3)$ என்பது முடிவிலுள்ள
!!{formula}~$\eq[s_1][s_3]$ ஆகவும் அமைகின்றன.
\end{ex}

% SEGMENT OLP-0110-S05
\begin{prob}
பின்வருவனவற்றுக்கு மூடிய !!{tableau}s தருக:
\begin{enumerate}
\item $\sFmla{\False}{\lforall[x][\lforall[y][((x = y \land !A(x))
      \lif !A(y))]]}$
\item $\sFmla{\False}{\lexists[x][(!A(x) \land
    \lforall[y][(!A(y) \lif y = x)])]}$,\\
  $\sFmla{\True}{\lexists[x][!A(x)] \land
    \lforall[y][\lforall[z][((!A(y) \land !A(z)) \lif y = z)]]}$
\end{enumerate}
\end{prob}

\end{document}
